\documentclass[10pt]{article}
\usepackage[margin=1in]{geometry}
\usepackage[T1]{fontenc}
\usepackage{lmodern,microtype,amsmath,amssymb,amsthm}
\theoremstyle{definition}
\newtheorem{assumption}{Assumption}
\begin{document}
\section*{Stochastic Generalized Monotone Root Finding}
We seek a zero of a continuous monotone operator $F:E\to E$ on a finite-dimensional real Hilbert space $E$. For known constants $L_0,L_1\geq0$, its variation along any segment is bounded by the endpoint distance times $L_0+L_1M$ whenever $M\in\mathbb{R}$ bounds the operator norm throughout that segment. A given point $x^0\in E$ lies within a known distance $R\geq0$ of a zero $x^\star$. Each oracle call returns an unbiased random vector estimating $F(x)$ at the query $x\in E$, using independent samples and with mean squared error at most a known $\sigma^2$, where $\sigma\geq0$. For $\varepsilon>0$ and $\delta\in(0,1)$, a measurable output $\bar{x}\in E$ must satisfy $\mathbb{P}(\|F(\bar{x})\|>\varepsilon)\leq\delta$, where $\mathbb{P}$ includes all algorithm and oracle randomness.

The root-finding problem is
\[
\text{find }x\in E\text{ such that }F(x)=0,
\]
where the norm $\|\cdot\|$ is induced by the inner product $\langle\cdot,\cdot\rangle$. The operator and its stochastic oracle satisfy the following conditions.

\begin{assumption}[Monotonicity and generalized Lipschitz continuity]
The operator $F$ is continuous, and for every $x,y\in E$,
\[
\langle F(x)-F(y),x-y\rangle\geq0.
\]
For every $x,y\in E$ and $M\in\mathbb{R}$, the segment bound
\[
\|F(x+\theta(y-x))\|\leq M
\qquad\text{for every }\theta\in[0,1]
\]
implies
\[
\|F(x)-F(y)\|\leq(L_0+L_1M)\|x-y\|.
\]
\end{assumption}

\begin{assumption}[Existence of a nearby zero]
There exists $x^\star\in E$ such that
\[
\begin{aligned}
F(x^\star)&=0,\\
\|x^0-x^\star\|&\leq R.
\end{aligned}
\]
\end{assumption}

\begin{assumption}[Unbiased stochastic operator]
Let $\Omega$ be a measurable sample space. There exist a probability measure $p$ on $\Omega$ and a jointly measurable map $g:E\times\Omega\to E$. For every $x\in E$, the map $g(x,\cdot)$ is Bochner integrable under $p$, the function $\omega\mapsto\|g(x,\omega)-F(x)\|^2$ is integrable under $p$, and
\[
\begin{aligned}
\int_\Omega g(x,\omega)\,p(d\omega)&=F(x),\\
\int_\Omega\|g(x,\omega)-F(x)\|^2\,p(d\omega)&\leq\sigma^2.
\end{aligned}
\]
\end{assumption}

\begin{assumption}[Independent sampling and oracle access]
There is a probability law $P$ on $\Omega^{\mathbb{N}}$ under which the coordinate samples $\xi_k$, $k\in\mathbb{N}$, are measurable, mutually independent, and each has distribution $p$. At call $k$, a query $x\in E$ returns the vector $g(x,\xi_k)\in E$ and costs one oracle call. Any internal algorithm randomness is independent of the entire sample sequence.
\end{assumption}

The operator is accessed through this oracle, and queries may be chosen adaptively anywhere in $E$. The public inputs are the space $E$, the measurable sample space $\Omega$, the initial point $x^0$, and all numerical parameters $(L_0,L_1,R,\sigma,\varepsilon,\delta)$ in the ranges stated above. The operator $F$, the oracle map $g$, the sampling laws $p$ and $P$, and the witnessing zero $x^\star$ are not supplied as public inputs. An externally specified budget $K\in\mathbb{N}$ may depend on the public parameter tuple and counts every oracle call.

After $K$ calls, the algorithm returns a measurable point $\bar{x}\in E$. Under the joint law $\mathbb{P}$ of its internal randomness and the oracle samples, the required residual accuracy is
\[
\mathbb{P}\bigl(\|F(\bar{x})\|>\varepsilon\bigr)\leq\delta.
\]
\end{document}
